problem:  
Let \( X = G/K \) be an irreducible symmetric space of noncompact type and rank one (for example, real hyperbolic space \( \mathbb{H}^n \) or complex hyperbolic space), and let \( Y = G'/K' \) be an irreducible symmetric space of noncompact type with \( \operatorname{rank}(Y) \ge 2 \).  

For any harmonic map \( f : X \to Y \) with **unbounded image**, define its **visual limit set**  
\[
\Lambda_f := \overline{f(X)} \cap \partial_\infty Y
\]
in the visual compactification of \( Y \).  

**Problem:**  
Must \( \Lambda_f \) have full Tits diameter, i.e. \( \mathrm{diam}_{Tits}(\Lambda_f) = \pi \)?  
Equivalently, can there exist a nonconstant harmonic map \( f : X \to Y \) whose image at infinity \( \Lambda_f \) is contained in a proper spherical subbuilding of the Tits boundary \( \partial_\infty Y \)?  

In other words, is there *analytic rank-collapse* at infinity: can harmonic maps from a negatively curved space be asymptotically confined to a lower-rank geometric sector of a higher-rank symmetric space without being totally geodesic?

---

Why is it a "good" problem:  
This question probes harmonic map rigidity from a **purely geometric–analytic viewpoint**, replacing the traditional topological/equivariant framework (as in Mostow or Corlette rigidity) by the intrinsic asymptotic geometry of harmonic maps.  It asks whether the analytic and curvature constraints of harmonicity can detect the algebraic structure of the Tits building at infinity of the target.  

Technically, it marries ideas from **harmonic analysis on symmetric spaces**, **asymptotic geometry**, and **building theory**: the Tits metric encodes the algebraic stratification of \( Y \) by rank, while harmonic map theory brings in curvature, convexity, and analytic inequalities. A positive answer would uncover a new form of rigidity where the “spread” of a harmonic map’s image at infinity must reflect the full rank of the target; a counterexample would demonstrate new analytic flexibility within higher-rank geometries.  

The statement is simple and geometric—asking whether a smooth harmonic map can avoid “seeing” parts of the boundary—but reaches into deep, unexplored terrain linking partial differential equations, differential geometry, and Lie theory. It thus exemplifies a **“narrow entrance, profound depth”** research problem: easy to grasp, but any resolution would reshape our understanding of how curvature and rank interact through harmonic mappings.